% ══════════════════════════════════════════════════════════════════════════
% §7  The "Don't Do This" Checklist
% ══════════════════════════════════════════════════════════════════════════
\section{The ``Don't Do This'' Checklist}
\label{sec:pitfalls}

Every bug below has crashed a real simulation at least once.
They are ordered from ``most silent'' to ``most spectacular.''

% ── Pitfall 1 ─────────────────────────────────────────────────────────
\subsection{Silent Failure: Ignoring Solver Status}

\begin{warningbox}
Just because a solver function \emph{returns} does not mean it
\emph{succeeded}.  Eigen solvers set an \texttt{info()} flag.
If you never check it, you will silently propagate garbage.
\end{warningbox}

\begin{lstlisting}
// BAD: no error check
Eigen::VectorXd x = solver.solve(b);
use(x);  // x might be nonsense

// GOOD: always check
Eigen::VectorXd x = solver.solve(b);
if (solver.info() != Eigen::Success) {
    std::cerr << "Solver failed: " << solver.info() << "\n";
    return EXIT_FAILURE;
}
\end{lstlisting}

% ── Pitfall 2 ─────────────────────────────────────────────────────────
\subsection{Comparing Floating-Point Values to Zero}

Floating-point arithmetic is \emph{approximate}.  After a chain of
operations, a value that should be exactly zero will instead be
$\sim 10^{-16}$.

\begin{lstlisting}
// BAD: will almost never trigger
if (determinant == 0.0) { /* singular */ }

// ALSO BAD: 1e-12 is an arbitrary "magic number".
// If matrix entries are ~1e-20, then det = 1e-15 is "huge".
// If entries are ~1e+10, then det = 1e-12 is effectively zero.
if (std::abs(determinant) < 1e-12) { /* fragile */ }

// GOOD: use a *relative* tolerance
double scale = A.norm();  // or max entry magnitude
if (std::abs(val) < epsilon * scale) { /* effectively zero */ }
\end{lstlisting}

\begin{tipbox}
Absolute tolerances like \texttt{1e-12} are \textbf{scale-dependent}
and will silently give wrong answers when your matrix entries are
very large or very small.  The most robust check for singularity is the
\emph{condition number} $\kappa(A) = \sigma_1 / \sigma_{\min}$,
which is scale-invariant.
If $\kappa > 1/\varepsilon_{\mathrm{mach}} \approx 10^{16}$,
the matrix is numerically singular regardless of its determinant.
\end{tipbox}

% ── Pitfall 3 ─────────────────────────────────────────────────────────
\subsection{Allocation Inside Hot Loops}

\texttt{Eigen::MatrixXd} dynamically allocates heap memory.
Allocating and freeing inside a tight loop is orders of magnitude
slower than the actual arithmetic.

\begin{lstlisting}
// BAD: allocates n x n on every iteration
while (running) {
    Eigen::MatrixXd A(n, n);   // malloc + free every loop
    fill(A);
    solve(A, b);
}

// GOOD: allocate once, reuse
Eigen::MatrixXd A(n, n);
while (running) {
    fill(A);
    solve(A, b);
}
\end{lstlisting}

For fixed-size small matrices ($n \le 4$), use \texttt{Eigen::Matrix3d}
or \texttt{Eigen::Matrix4d}---they live on the stack and cost zero
allocation overhead.

% ── Pitfall 4 ─────────────────────────────────────────────────────────
\subsection{Skipping the Residual Check}

This is the single most important engineering habit in numerical
computing.  After every solve, compute:
\[
    \text{relative residual} = \frac{\|Ax - b\|}{\|b\|}.
\]

\begin{lstlisting}
Eigen::VectorXd x = solver.solve(b);
double rel_res = (A * x - b).norm() / b.norm();
if (rel_res > 1e-8) {
    std::cerr << "WARNING: residual = " << rel_res
              << " -- solution may be inaccurate\n";
}
\end{lstlisting}

\begin{warningbox}
A small residual does \emph{not} guarantee a small error in~$x$.
If $\kappa(A)$ is large, even a tiny residual can correspond to
a wildly wrong solution.  Always pair the residual check with
a condition number estimate when possible.
\end{warningbox}

% ── Pitfall 5 ─────────────────────────────────────────────────────────
\subsection{Aliasing in Eigen Expressions}

Eigen uses lazy evaluation.  The classic ``silent killer'' is
\texttt{A = A * A}: the same matrix appears on both sides of the
assignment, and without intervention Eigen would read and write
the same memory simultaneously, producing garbage.

In modern Eigen ($\ge 3.3$), the library detects \emph{some}
aliasing patterns automatically---for example,
\texttt{A = A * B} (where \texttt{B} is a \emph{different} matrix)
is handled correctly.  However, \texttt{A = A * A} remains the
classic case that requires explicit intervention.

\begin{lstlisting}
// DANGEROUS: same matrix on both sides
A = A * A;   // Eigen 3.3+ may catch this, but don't rely on it

// SAFE: .eval() is cheap insurance for any self-referencing expression
A = (A * A).eval();
// or equivalently:
Eigen::MatrixXd temp = A * A;
A = temp;
\end{lstlisting}

\begin{tipbox}
Rule of thumb: if the \textbf{same matrix name appears on both sides}
of the equals sign, add \texttt{.eval()}.  The cost is one temporary
allocation---trivial compared to the hours you will spend debugging
silent corruption.
\end{tipbox}

\subsection{Complete Pitfall Demonstration}

\lstinputlisting[caption={All five pitfalls demonstrated and fixed.},
                 label={lst:pitfalls}]
                {code/pitfalls_demo.cpp}
